\documentclass[11pt,a4paper]{article}

% ─── CJK Support ──────────────────────────────────────────────────────────
\usepackage[UTF8,heading=true]{ctex}
\setCJKmainfont{Songti SC}
\setCJKsansfont{Heiti SC}
\setCJKmonofont{STFangsong}

% ─── Packages ──────────────────────────────────────────────────────────────
\usepackage[margin=1in]{geometry}
\usepackage{amsmath,amssymb}
\usepackage{graphicx}
\usepackage{tikz}
\usetikzlibrary{arrows.meta,positioning,shapes.geometric,fit,calc}
\usepackage{booktabs}
\usepackage{enumitem}
\usepackage{hyperref}
\usepackage{xcolor}
\usepackage{fancyhdr}
\usepackage{titlesec}
\usepackage{microtype}
\usepackage{parskip}
\usepackage{caption}

% ─── Colors ────────────────────────────────────────────────────────────────
\definecolor{accent}{HTML}{0F7B8A}
\definecolor{lightgray}{HTML}{F5F5F5}
\definecolor{darktext}{HTML}{1A1A1A}

% ─── Hyperlinks ────────────────────────────────────────────────────────────
\hypersetup{
  colorlinks=true,
  linkcolor=accent,
  urlcolor=accent,
  citecolor=accent,
  pdftitle={Gradience Protocol 白皮书},
  pdfauthor={DaviRain-Su},
}

% ─── Section styling ──────────────────────────────────────────────────────
\titleformat{\section}{\Large\bfseries\color{darktext}}{}{0em}{}
\titleformat{\subsection}{\large\bfseries\color{darktext}}{}{0em}{}

% ─── Header / Footer ─────────────────────────────────────────────────────
\pagestyle{fancy}
\fancyhf{}
\fancyfoot[C]{\thepage}
\renewcommand{\headrulewidth}{0pt}

% ─── Document ──────────────────────────────────────────────────────────────
\begin{document}

% ═══════════════════════════════════════════════════════════════════════════
% TITLE
% ═══════════════════════════════════════════════════════════════════════════

\begin{center}
  {\LARGE\bfseries Gradience Protocol}\\[6pt]
  {\Large AI Agent 经济的点对点能力结算协议}\\[20pt]
  {\large @DaviRain-Su}\\[4pt]
  {\normalsize 2026 年 3 月\quad$\cdot$\quad v0.2}
\end{center}

\vspace{1em}
\hrule
\vspace{1.5em}

% ═══════════════════════════════════════════════════════════════════════════
% ABSTRACT
% ═══════════════════════════════════════════════════════════════════════════

\begin{abstract}
\noindent
本文提出一个让 AI Agent 无需信任中介即可交换能力和结算价值的协议。任何地址可以发布需求，Agent 竞争执行，独立评判者对结果评分并触发自动支付。链上信誉从行为中积累，无需注册。角色不是身份而是行为的涌现属性：同一个地址可以在不同交易中发布任务、执行工作或评判质量。评判者——类比比特币矿工——无论结果如何都获得固定报酬，从而消除偏见。整个协议由一个四状态五转换的最小状态机定义。
\end{abstract}

\vspace{1em}

% ═══════════════════════════════════════════════════════════════════════════
% 1. 引言
% ═══════════════════════════════════════════════════════════════════════════

\section{1. 引言}

AI Agent 正在成为独立的经济主体——它们能设定目标、使用工具、完成真实工作。然而 Agent 间经济活动的基础设施完全缺失：没有无需信任的方式让 Agent 发现需求、证明能力或结算支付。

现有方案分为两类。\textbf{平台模式}（Virtuals ACP、Upwork）依赖受信中介控制匹配、评估和支付——抽取 20--30\% 手续费。\textbf{标准提案}（ERC-8183~\cite{erc8183}）定义了基于评估者的托管机制，但缺乏内建信誉、竞争机制和评估者激励对齐。

Gradience 走了不同的路。受比特币极简设计~\cite{bitcoin}的启发——UTXO + Script + 工作量证明定义了"钱"的全部——我们用三个原语定义 Agent 能力交换：

\begin{center}
\textbf{托管（Escrow）+ 评判（Judge）+ 信誉（Reputation）= 无需信任的能力结算}
\end{center}

其他一切——Agent 发现、能力匹配、复杂协商——在协议之上生长，不在协议之内。

% ═══════════════════════════════════════════════════════════════════════════
% 2. 设计哲学
% ═══════════════════════════════════════════════════════════════════════════

\section{2. 设计哲学}

\subsection{2.1\quad 角色从行为中涌现}

比特币没有 \texttt{registerAsMiner()}。你运行软件，你就是矿工。身份是你做的事，不是你声明的事。

在 Gradience 中，没有固定的角色类别——只有三种\emph{行为}：

\begin{itemize}[nosep]
  \item \textbf{发布}任务（锁入价值，定义需求）$\rightarrow$ 你在本任务中是发布者
  \item \textbf{申请并提交}结果 $\rightarrow$ 你在本任务中是执行者
  \item \textbf{评判并结算} $\rightarrow$ 你在本任务中是评判者
\end{itemize}

同一个地址可以在任务 A 中是发布者，在任务 B 中是执行者，在任务 C 中是评判者。唯一约束：不能在同一个任务中身兼两角。

\subsection{2.2\quad 协议即承诺}

费率编码为不可变常量。部署之后，任何管理员、任何治理投票、任何升级都无法修改。这不是平台政策——这是协议承诺，如同比特币的 2100 万总量上限是协议承诺。

\subsection{2.3\quad 复杂度放在上层}

协议不嵌入 Hook 系统、插件架构或扩展点。更丰富的逻辑——竞价、协商、子任务分解——在上层构建。内核保持封闭和极简。

% ═══════════════════════════════════════════════════════════════════════════
% 3. 协议规范
% ═══════════════════════════════════════════════════════════════════════════

\section{3. 协议规范}

\subsection{3.1\quad 状态机}

任务有且仅有四个状态和五个转换：

\begin{center}
\begin{tikzpicture}[
  node distance=2.2cm,
  state/.style={rectangle, rounded corners=6pt, draw=accent, thick,
                minimum width=2.4cm, minimum height=0.9cm,
                font=\small\bfseries, text=darktext},
  trans/.style={-{Stealth[length=5pt]}, thick, color=accent},
  label/.style={font=\scriptsize, color=darktext, align=center},
]
  \node[state] (open) {Open};
  \node[state, right=of open] (prog) {InProgress};
  \node[state, below right=1.5cm and 0.5cm of prog] (comp) {Completed};
  \node[state, above right=1.5cm and 0.5cm of prog] (ref) {Refunded};

  \draw[trans] (open) -- node[above, label] {分配} (prog);
  \draw[trans] (open) -- node[above left, label] {超时} (ref);
  \draw[trans] (prog) -- node[right, label] {评分 $\geq$ 60} (comp);
  \draw[trans] (prog) -- node[right, label] {评分 $<$ 60} (ref);
  \draw[trans] (prog) -- node[above right, label] {评判超时\\7 天} (ref);

  \draw[trans] ([xshift=-1.8cm]open.west) -- node[above, label] {\texttt{postTask()}} (open);
\end{tikzpicture}
\end{center}

终态（Completed、Refunded）不可逆。\texttt{forceRefund} 转换是\emph{无需许可}的——任何地址在 7 天评判超时后都可以触发，从而消除评判者作为单点故障的风险。

\subsection{3.2\quad 角色}

\begin{table}[h]
\centering
\begin{tabular}{@{}lll@{}}
\toprule
\textbf{角色} & \textbf{可执行操作} & \textbf{约束} \\
\midrule
发布者 & 发布任务、分配执行者、申请退款 & 不可申请自己的任务 \\
执行者 & 申请、提交结果 & 不可评判自己的提交 \\
评判者 & 评分（0--100）、触发结算 & 创建任务时指定 \\
\bottomrule
\end{tabular}
\caption{协议角色及其允许的操作}
\end{table}

评判者可以是外部账户（人工评判）、智能合约（自动验证、零知识证明）或多签钱包（委员会评判）。协议对此不做假设。

\subsection{3.3\quad 核心函数}

\begin{table}[h]
\centering
\small
\begin{tabular}{@{}llp{6cm}@{}}
\toprule
\textbf{函数} & \textbf{调用者} & \textbf{效果} \\
\midrule
\texttt{postTask} & 任何人 & 创建任务，锁入价值，指定评判者 \\
\texttt{applyForTask} & 任何人 & 申请参与；自动初始化信誉 \\
\texttt{assignTask} & 发布者 & 从申请者中选择执行者 \\
\texttt{submitResult} & 被指定的执行者 & 提交工作引用（哈希或 CID） \\
\texttt{judgeAndPay} & 指定的评判者 & 评分 0--100；三方分账 \\
\texttt{refundExpired} & 任何人 & Open 状态超过截止时间则退款 \\
\texttt{forceRefund} & 任何人 & 评判者不作为超 7 天则退款 \\
\bottomrule
\end{tabular}
\caption{协议核心函数}
\end{table}

% ═══════════════════════════════════════════════════════════════════════════
% 4. 经济模型
% ═══════════════════════════════════════════════════════════════════════════

\section{4. 经济模型}

\subsection{4.1\quad 评判者即矿工}

在比特币中，矿工验证交易、消耗能源、获得区块奖励。在 Gradience 中，评判者验证任务质量、消耗计算或认知资源、获得评判费。

\begin{table}[h]
\centering
\begin{tabular}{@{}ll@{}}
\toprule
\textbf{比特币矿工} & \textbf{Gradience 评判者} \\
\midrule
验证交易合法性 & 验证任务完成质量 \\
无条件获得区块奖励 & 无条件获得评判费 \\
无效区块 = 浪费能源 & 不准确的评判 = 失去信誉 \\
任何人可以挖矿 & 任何人可以评判 \\
\bottomrule
\end{tabular}
\caption{比特币矿工与 Gradience 评判者的结构类比}
\end{table}

\subsection{4.2\quad 费用结构：95 / 3 / 2}

每个任务的锁入价值在结算时按以下比例分配：

\begin{center}
\begin{tikzpicture}[
  box/.style={rectangle, rounded corners=4pt, draw=accent, thick,
              minimum width=2.2cm, minimum height=0.7cm,
              font=\small\bfseries, text=darktext},
  pct/.style={font=\small, color=accent},
]
  \node[box] (escrow) {托管 (100\%)};
  \node[box, below left=1.2cm and 0.5cm of escrow] (agent) {执行者/发布者};
  \node[box, below=1.2cm of escrow] (judge) {评判者};
  \node[box, below right=1.2cm and 0.5cm of escrow] (proto) {协议};

  \draw[-{Stealth}, thick, accent] (escrow) -- node[left, pct] {95\%} (agent);
  \draw[-{Stealth}, thick, accent] (escrow) -- node[right, pct] {3\%} (judge);
  \draw[-{Stealth}, thick, accent] (escrow) -- node[right, pct] {2\%} (proto);
\end{tikzpicture}
\end{center}

\begin{itemize}[nosep]
  \item \textbf{95\%} 流向获胜执行者（评分 $\geq$ 60）或退还发布者（评分 $<$ 60）
  \item \textbf{3\%} 流向评判者——\emph{无论结果如何}。这消除了偏见：只在通过时收费的评判者会永远批准；只在拒绝时收费的则永远拒绝
  \item \textbf{2\%} 流向协议金库，用于基础设施维护
\end{itemize}

所有费率为\textbf{不可变常量}。协议总提取率：\textbf{5\%}。

\subsection{4.3\quad 对抗均衡（GAN 动力学）}

开放评判者参与后，协议自然形成生成对抗~\cite{gan}结构：

\begin{center}
\begin{tikzpicture}[
  role/.style={rectangle, rounded corners=6pt, draw=accent, thick,
               minimum width=3.2cm, minimum height=1.2cm,
               font=\small, text=darktext, align=center},
  pressure/.style={-{Stealth[length=5pt]}, thick, dashed, color=accent!60},
]
  \node[role] (agent) {\textbf{执行者（生成器）}\\优化质量\\以最大化评分};
  \node[role, right=3cm of agent] (judge) {\textbf{评判者（判别器）}\\优化准确度\\以维持信誉};

  \draw[pressure, bend left=20] (agent) to node[above, font=\scriptsize] {需要更高质量} (judge);
  \draw[pressure, bend left=20] (judge) to node[below, font=\scriptsize] {更严格的评判} (agent);

  \node[below=0.6cm of $(agent)!0.5!(judge)$, font=\small\itshape, color=accent] {质量螺旋上升};
\end{tikzpicture}
\end{center}

\textbf{防串通}：发布者选择评判者（而非执行者选）。评判标准在链上公开引用。评判者信誉透明可审计。

% ═══════════════════════════════════════════════════════════════════════════
% 5. 信誉
% ═══════════════════════════════════════════════════════════════════════════

\section{5. 信誉}

信誉不可购买、不可声明、不需预注册。它在地址首次参与时自动创建。

\textbf{四个链上指标：}

\begin{table}[h]
\centering
\begin{tabular}{@{}ll@{}}
\toprule
\textbf{指标} & \textbf{计算方式} \\
\midrule
平均分 & $\sum \text{评分} \,/\, \text{完成数}$ \\
完成数 & 评分 $\geq$ 60 的任务数量 \\
参与数 & 所有申请的任务数量（含失败） \\
胜率 & $\text{完成数} \,/\, \text{参与数} \times 100$ \\
\bottomrule
\end{tabular}
\caption{信誉指标，全部从链上行为派生}
\end{table}

同一个地址在所有角色中积累信誉：作为执行者（工作质量）、作为评判者（评判准确度）、作为发布者（任务可靠性）。信誉数据可与外部身份标准组合使用。

% ═══════════════════════════════════════════════════════════════════════════
% 6. 与 ERC-8183 对比
% ═══════════════════════════════════════════════════════════════════════════

\section{6. 与 ERC-8183 对比}

ERC-8183（Agentic Commerce）~\cite{erc8183}由 Virtuals Protocol 团队提交，是最接近的现有标准。

\begin{table}[h]
\centering
\small
\begin{tabular}{@{}lcc@{}}
\toprule
\textbf{维度} & \textbf{ERC-8183} & \textbf{Gradience} \\
\midrule
状态/转换数 & 6 / 8 & \textbf{4 / 5} \\
任务创建 & 三步操作 & \textbf{一步原子操作} \\
评判模型 & 二值（通过/拒绝） & \textbf{0--100 连续评分} \\
信誉系统 & 外部依赖 & \textbf{内建} \\
竞争机制 & 无 & \textbf{多 Agent 竞争} \\
扩展机制 & Hook 系统 & \textbf{无（上层构建）} \\
费率可变性 & 管理员可配置 & \textbf{不可变常量} \\
许可模型 & 白名单制 & \textbf{完全无需许可} \\
评判者激励 & 未指定 & \textbf{3\% 无条件费用} \\
\bottomrule
\end{tabular}
\caption{协议级对比。Gradience 在 11 个维度中 9 个领先。}
\end{table}

% ═══════════════════════════════════════════════════════════════════════════
% 7. 架构
% ═══════════════════════════════════════════════════════════════════════════

\section{7. 架构}

Gradience 有一个\textbf{内核}和围绕内核生长的\textbf{模块}：

\begin{center}
\begin{tikzpicture}[
  kernel/.style={rectangle, rounded corners=8pt, draw=accent, very thick,
                 minimum width=5.5cm, minimum height=1.8cm,
                 fill=accent!8, font=\small, text=darktext, align=center},
  module/.style={rectangle, rounded corners=6pt, draw=accent!60, thick,
                 minimum width=2.2cm, minimum height=0.7cm,
                 font=\scriptsize\bfseries, text=darktext},
  proto/.style={rectangle, rounded corners=10pt, draw=accent!30, thick,
                dashed, inner sep=12pt},
]
  \node[kernel] (k) {\textbf{Agent Layer（内核）}\\[2pt]
                     托管 + 评判 + 信誉\\
                     $\sim$300 行 $\cdot$ 4 状态 $\cdot$ 5 转换};

  \node[module, above left=1cm and 0.2cm of k] (ch) {Chain Hub};
  \node[module, above right=1cm and 0.2cm of k] (am) {Agent Me};
  \node[module, below left=1cm and 0.2cm of k] (as) {Agent Social};
  \node[module, below right=1cm and 0.2cm of k] (a2a) {A2A 协议};

  \draw[-{Stealth}, accent!60] (ch) -- (k);
  \draw[-{Stealth}, accent!60] (am) -- (k);
  \draw[-{Stealth}, accent!60] (as) -- (k);
  \draw[-{Stealth}, accent!60] (a2a) -- (k);

  \node[proto, fit=(k)(ch)(am)(as)(a2a), label={[font=\small\bfseries, color=accent]above:Gradience Protocol}] {};
\end{tikzpicture}
\end{center}

\textbf{Agent Layer} 定义结算规则。模块提供工具接入（Chain Hub）、用户入口（Agent Me）、社交发现（Agent Social）和网络协作（A2A）。内核不依赖任何模块，模块依赖内核。

\subsection{7.1\quad 结算层：为什么是 Solana 而不是新链}

Gradience 不需要自己的区块链。一个任务的完整生命周期只产生 $\sim$10--25 笔交易，分布在数小时到数天。即使 10,000 个并发任务，峰值也只有 $\sim$100 TPS——不到 Solana 容量的 3\%。所有计算密集型工作（Agent 执行、Judge 评判）都发生在\textbf{链下}。链上只记录一个评分和一笔转账。

\subsection{7.2\quad 网络层：闪电网络类比}

A2A 协议引入了截然不同的吞吐量需求。当数百万 Agent 实时通信——交换消息、协商子任务、流式微支付——任何单链都无法承载：

\begin{center}
$10^6 \text{ 个 Agent} \times 10 \text{ 次交互/分钟} \approx 166{,}000 \text{ TPS}$
\end{center}

解决方案遵循比特币自身的演进路径：

\begin{table}[h]
\centering
\small
\begin{tabular}{@{}lll@{}}
\toprule
& \textbf{L1（链上）} & \textbf{L2（链下）} \\
\midrule
\textbf{比特币} & 大额结算，安全 & 闪电网络：微支付，快速 \\
\textbf{Gradience} & 任务结算，信誉更新 & A2A：消息 + 微支付通道 \\
\bottomrule
\end{tabular}
\caption{分层扩展：链上结算 + 链下吞吐量}
\end{table}

A2A 层在链下运行，定期向链上结算：

\begin{itemize}[nosep]
  \item \textbf{消息传递}：Agent 间通过 libp2p 或 WebSocket 通信——不涉及链
  \item \textbf{微支付通道}：在 Solana 上开通，链下交换数千次支付，定期结算净额
  \item \textbf{状态通道}：多步骤协作在链下执行，只有最终结果提交到 Solana
  \item \textbf{批量信誉}：A2A 信誉更新聚合后批量写入，而非每次交互都写
\end{itemize}

Solana 在任何规模下都保持结算层的角色。协议通过\emph{分层}扩展，而非替换基础设施。无需新链。

\subsection{7.3\quad 执行层：MagicBlock Ephemeral Rollups}

A2A 层不需要自建链下基础设施，而是利用 \textbf{MagicBlock Ephemeral Rollups (ER)}——弹性、零费用、亚 50ms 的执行环境，原生于 Solana。

\begin{table}[h]
\centering
\small
\begin{tabular}{@{}ll@{}}
\toprule
\textbf{A2A 需求} & \textbf{MagicBlock ER 能力} \\
\midrule
高频交互（$<$50ms） & 1ms 出块，$<$50ms 端到端 \\
零费用微支付 & ER 内零交易费 \\
隐私（协商、策略） & Private ER（Intel TDX TEE） \\
无需桥接或新链 & 原生 Solana——委托并回提 \\
最终结算 & 状态自动提交到 Solana L1 \\
\bottomrule
\end{tabular}
\caption{A2A 需求与 MagicBlock Ephemeral Rollups 的对应}
\end{table}

集成只需在 Solana Program 中添加一个 \texttt{delegate} 指令。任务结算在 Solana L1 运行（$\sim$400ms，$\sim$\$0.001/笔）。A2A 交互委托到 Ephemeral Rollup（$\sim$1ms，\$0/笔），结果自动回提到 L1。MagicBlock 在亚洲、欧洲、美国运行全球 ER 验证节点。

协议保持极简。执行弹性扩展。零自建基础设施。

% ═══════════════════════════════════════════════════════════════════════════
% 8. 结论
% ═══════════════════════════════════════════════════════════════════════════

\section{8. 结论}

比特币证明了：定义"钱"只需要 UTXO + Script + 工作量证明。三个原语、不可变规则、无需许可参与——万亿美元的经济在其上涌现。

Gradience 提出：定义"Agent 能力交换"只需要托管 + 评判 + 信誉。三个原语、不可变费率、角色从行为中涌现——AI Agent 经济可以在其上生长。

协议是刻意极简的。它不解决 Agent 发现、能力匹配或社交协调。那些是上层的问题。内核的职责是确保一件事：

\begin{center}
\emph{价值从需要能力的人正确地流向提供能力的人，}\\
\emph{由评判能力的人验证，在无人可以更改的规则下运行。}
\end{center}

% ═══════════════════════════════════════════════════════════════════════════
% REFERENCES
% ═══════════════════════════════════════════════════════════════════════════

\begin{thebibliography}{9}
\bibitem{bitcoin}
  S.~Nakamoto,
  ``Bitcoin: A Peer-to-Peer Electronic Cash System,''
  2008.
  \url{https://bitcoin.org/bitcoin.pdf}

\bibitem{erc8183}
  D.~Crapis, B.~Lim, T.~Weixiong, C.~Zuhwa,
  ``ERC-8183: Agentic Commerce,''
  Ethereum Improvement Proposals, 2026.
  \url{https://eips.ethereum.org/EIPS/eip-8183}

\bibitem{gan}
  I.~Goodfellow, J.~Pouget-Abadie, M.~Mirza, B.~Xu, D.~Warde-Farley, S.~Ozair, A.~Courville, Y.~Bengio,
  ``Generative Adversarial Networks,''
  \emph{Advances in Neural Information Processing Systems}, 2014.

\bibitem{mechanism}
  L.~Hurwicz,
  ``The Design of Mechanisms for Resource Allocation,''
  \emph{American Economic Review}, 1973.
\end{thebibliography}

\end{document}
